\setlength{\absparsep}{18pt} % ajusta o espaçamento dos parágrafos do resumo
\begin{resumo}

Modelos de aprendizado de máquina treinados com microdados educacionais podem reproduzir e amplificar desigualdades estruturais sem que as métricas convencionais de desempenho revelem o problema. Este trabalho apresenta o AuditEdu-XAI, uma ferramenta de software de auditoria algorítmica que detecta, quantifica e explica vieses em modelos preditivos de desempenho no Exame Nacional do Ensino Médio (ENEM). Desenvolvida segundo a metodologia de \textit{Design Science Research}, a ferramenta adota uma arquitetura desacoplada: um pipeline em lote processa os microdados das edições de 2019, 2020 e 2021, treina classificadores XGBoost e calcula, exclusivamente sobre o conjunto de teste, o Impacto Disparate, a Paridade Demográfica, as taxas de verdadeiros e falsos positivos (\textit{Equal Opportunity} e \textit{Equalized Odds}) e explicações TreeSHAP. Os resultados são persistidos em arquivos Parquet estáticos e consumidos por uma interface Streamlit com tempo de resposta inferior a 2 segundos em validação local. A interface está publicada gratuitamente no Streamlit Community Cloud e empacotada em contêiner Docker validado por integração contínua. Aplicada a modelos com AUC-ROC em torno de 0,75, a auditoria revelou que todos os grupos não brancos e os candidatos com renda familiar de até 2,5 salários mínimos violam a regra dos quatro quintos nas três edições, um resultado que se mantém nos intervalos de confiança por \textit{bootstrap}. Revelou também que o modelo amplifica a disparidade de raça/cor presente nos dados em 14 de 15 combinações de grupo e edição, 13 delas com intervalo de confiança inteiramente negativo, e que a ausência de computador na residência é o atributo mais influente nas predições, com peso por candidato crescente entre 2019 e 2021, em coincidência temporal com a pandemia. Para validar o artefato, foi conduzido um experimento de injeção sintética de viés com controle negativo: ao rebaixar 50\% dos rótulos positivos de candidatos pardos, o Impacto Disparate do grupo caiu para valores próximos de zero, enquanto a AUC-ROC variou, em média, no máximo 0,032. O resultado indica que a ferramenta detecta e localiza vieses aos quais as métricas convencionais de desempenho são pouco sensíveis.

\vspace{\onelineskip}

\noindent
\textbf{Palavras-chave}: justiça algorítmica; inteligência artificial explicável; impacto disparate; SHAP; ENEM.
\end{resumo}
